\section{Деревья и их основные свойства}
\label{sec:trees-properties}

\subsection{Определение дерева и эквивалентные свойства}

Лесом называется неориентированный граф без циклов. Связный лес называется
\emph{деревом}. Пусть $G=(V,E)$, $|V|=n$. Для конечного простого неориентированного
графа эквивалентны следующие утверждения:
\begin{enumerate}
    \item $G$ --- дерево;
    \item $G$ связен и имеет $n-1$ ребро;
    \item $G$ ацикличен и имеет $n-1$ ребро;
    \item любые две различные вершины соединены единственным простым путём;
    \item $G$ связен, но удаление любого ребра нарушает связность;
    \item $G$ ацикличен, но добавление любого нового ребра создаёт ровно один цикл.
\end{enumerate}

Это один из главных теоретических результатов о деревьях. Схема доказательства
строится по кругу. Из связности и отсутствия циклов следует единственность пути:
два разных пути между теми же вершинами образовали бы цикл. Если удалить ребро
дерева, единственный путь между его концами исчезает, поэтому граф распадается.
Последовательно удаляя листья, получают число рёбер $n-1$. Обратно, связный граф
с $n-1$ рёбрами не может содержать цикл: удаление ребра цикла сохранило бы
связность и дало связный граф с меньшим, чем $n-1$, числом рёбер.

Для леса с $n$ вершинами и $k$ компонентами связности
\[
    m=n-k.
\]
Действительно, каждая компонента является деревом; суммирование равенств
$m_i=n_i-1$ даёт $m=n-k$.

\subsection{Листья и степени вершин}

Вершина степени $1$ называется листом, или висячей вершиной. Всякое дерево с
$n\ge2$ имеет не менее двух листьев. Один из способов доказательства --- взять
максимальный простой путь: его концы не могут иметь соседей вне пути, иначе путь
можно продолжить; и не могут иметь дополнительный сосед внутри пути, иначе
возникнет цикл.

По лемме о рукопожатиях
\[
    \sum_{v\in V}d(v)=2(n-1).
\]
Если $n_k$ --- число вершин степени $k$, то
\[
    \sum_{k\ge1}n_k=n,
    \qquad
    \sum_{k\ge1}k n_k=2(n-1).
\]
Вычитая удвоенное первое равенство из второго, получаем
\[
    n_1=2+\sum_{k\ge3}(k-2)n_k.
\]
Эта формула показывает, что вершины высокой степени ``порождают'' дополнительные
листья.

\subsection{Корневые деревья}

Если выделена вершина $r$, дерево называется корневым. Единственный путь из $r$
в вершину $v$ определяет родителя $p(v)$ и цепочку предков. Соседи, для которых
$v$ является родителем, называются детьми. Глубина вершины
\[
    depth(v)=d(r,v),
\]
а высота дерева --- максимальная глубина.

Поддерево вершины $v$ состоит из $v$ и всех её потомков. Если направить все рёбра
от корня, получим арборесценцию:
\[
 d_{\rm in}(r)=0,
 \qquad
 d_{\rm in}(v)=1\quad(v\ne r),
\]
и каждая вершина достижима из корня ровно одним ориентированным путём.

\subsection{Бинарные деревья}

Упорядоченным называется корневое дерево, в котором для детей каждой вершины задан
порядок. Бинарное дерево имеет не более двух различаемых детей --- левого и
правого.

Терминология разных книг отличается, поэтому полезно явно задавать определение.
Если каждая внутренняя вершина имеет ровно двух детей, такое дерево называют
строгим (часто также полным) бинарным. Для него
\[
    L=I+1,
\]
где $L$ --- число листьев, $I$ --- число внутренних вершин. Доказательство:
число рёбер равно $2I$, но одновременно оно равно $(I+L)-1$.

Если высота бинарного дерева равна $h$ и корень имеет уровень $0$, то
\[
    n\le 1+2+\cdots+2^h=2^{h+1}-1.
\]
Отсюда для дерева с $n$ вершинами
\[
    h\ge \lceil \log_2(n+1)\rceil-1.
\]
Именно поэтому сбалансированные бинарные деревья поиска дают логарифмическую
высоту.

\subsection{Центр и диаметр дерева}

Диаметр дерева --- длина максимального простого пути. Центр --- множество вершин
минимального эксцентриситета
\[
 e(v)=\max_u d(u,v).
\]
\textbf{Теорема.} Центр дерева состоит либо из одной вершины, либо из двух
смежных вершин. Его можно получить, одновременно удаляя все листья до тех пор,
пока не останется одна или две вершины.

Диаметр можно найти двумя обходами:
\begin{enumerate}
 \item из произвольной вершины найти самую далёкую вершину $u$;
 \item из $u$ найти самую далёкую вершину $v$.
\end{enumerate}
Тогда путь $u$--$v$ является диаметром. Для невзвешенного дерева оба шага
выполняются BFS/DFS за $O(n)$.

\subsection{Остовные деревья}

Остовное дерево связного графа $G$ --- подграф, содержащий все вершины $G$ и
являющийся деревом. Связный граф имеет остовное дерево: рёбра дерева BFS или DFS,
по которым впервые достигаются новые вершины, образуют такой остов.

Для взвешенного графа возникает задача \emph{минимального остовного дерева}
(MST). Центральное свойство --- \textbf{свойство разреза}: если некоторый разрез
графа пересекает ребро минимального веса (и выбор однозначен), то это ребро можно
включить в некоторое MST. Именно это обосновывает алгоритмы Краскала и Прима.

\textbf{Алгоритм Краскала.} Сортировать рёбра по весу и добавлять их в порядке
возрастания, если они соединяют разные компоненты. Обычно компоненты поддерживают
структурой DSU. Сложность
\[
    O(m\log m)
\]
из-за сортировки.

\textbf{Алгоритм Прима.} Начать с одной вершины и каждый раз добавлять самое
лёгкое ребро, соединяющее уже построенную часть с внешней вершиной. С бинарной
кучей и списками смежности типичная сложность
\[
    O(m\log n).
\]
Если все веса рёбер попарно различны, MST единственно.

\subsection{Код Прюфера и формула Кэли}

Помеченному дереву на вершинах $\{1,\ldots,n\}$ ставится в соответствие код
Прюфера длины $n-2$: на каждом шаге удаляют лист с наименьшей меткой и записывают
метку его соседа. Процесс заканчивается, когда остаются две вершины.

Код Прюфера однозначно определяет исходное дерево. Число появлений вершины $v$ в
коде равно
\[
    d(v)-1.
\]
Поскольку любая последовательность длины $n-2$ над алфавитом из $n$ меток
является кодом некоторого дерева, получаем \textbf{формулу Кэли}
\[
    N_n=n^{n-2}
\]
для числа помеченных деревьев на $n$ вершинах.

\subsection{Что важно сказать на экзамене}

Нужно начать с эквивалентных характеристик дерева: связность + ацикличность,
$n-1$ ребро, единственность пути. Далее --- листья, корневое дерево и бинарные
деревья. Для связи с алгоритмами полезно рассказать об остовных деревьях и
свойстве разреза, которое обосновывает Краскала и Прима. В качестве сильного
дополнения можно привести код Прюфера и формулу Кэли $n^{n-2}$.

\newpage
